% ============================================================================
% Chapter 12: Software - Python
% Book source: part1-linear-programming/ch03-software/software-python-book1.tex
% Exercise numbering synced to \begin{ex} order as of 2026-07-13.
% All code executed with Python 3.10, PuLP 3.3.2, bundled CBC solver.
% ============================================================================

\section{Chapter 12: Software -- Python}

This chapter has nine exercises, all built around the two running PuLP
examples of the chapter's PuLP section (Section 12.12): the product mix LP
and the complete transportation LP. The warm-ups (12.1--12.3) modify data in the worked
examples; the core problems (12.4--12.6) build a model from scratch, extract
duals and slacks, and run an $\varepsilon$-constraint sweep; the concept
problems (12.7--12.8) diagnose the two classic PuLP modeling bugs (a missing
constraint and a wrong optimization sense); and the challenge problem (12.9)
traces optimal profit as a function of a right-hand side. Every solution below
was produced by actually running the code shown (PuLP 3.3.2 with the bundled
CBC solver); printed output is reproduced verbatim. The book's Selected
Solutions for 12.2, 12.5, and 12.9 were re-executed and confirmed; they are
reused and expanded here.

% ----------------------------------------------------------------------------
\exsol{12.1}{A New Profit Margin}{1}

Only the objective line changes: the profit coefficient of $X_1$ becomes 5.

\begin{verbatim}
from pulp import *

prob = LpProblem("Product_Mix", LpMaximize)
x1 = LpVariable("X1", lowBound=0)
x2 = LpVariable("X2", lowBound=0)

prob += 5*x1 + 2*x2, "Total_Profit"     # was 3*x1 + 2*x2
prob += 10*x1 + 5*x2 <= 300, "Material"
prob += 4*x1 + 4*x2 <= 160, "Labor"
prob += 2*x1 + 6*x2 <= 180, "Machine"

prob.solve(PULP_CBC_CMD(msg=0))
print(f"Status: {LpStatus[prob.status]}")
print(f"Maximum Profit: ${value(prob.objective)}")
print(f"X1 = {x1.varValue}")
print(f"X2 = {x2.varValue}")
for name, c in prob.constraints.items():
    print(name, "slack =", c.slack)
\end{verbatim}

Actual output:

\begin{verbatim}
Status: Optimal
Maximum Profit: $150.0
X1 = 30.0
X2 = 0.0
Material slack = -0.0
Labor slack = 40.0
Machine slack = 120.0
\end{verbatim}

The new plan is $(X_1, X_2) = (30, 0)$ with profit $\$150$. Product 1 is now
so profitable per unit of the scarce material (profit $5$ per $10$ units of
material, versus $2$ per $5$ for Product 2, and material is the first resource
to run out) that the company specializes completely. Only the
\textbf{Material} constraint is binding ($10 \cdot 30 + 5 \cdot 0 = 300$);
Labor has slack $40$ and Machine has slack $120$ (the nonnegativity bound
$X_2 \geq 0$ is also active). As promised, the answer is no longer $(20,20)$.

% ----------------------------------------------------------------------------
\exsol{12.2}{From Math to PuLP}{1}

The book's Selected Solution is correct as printed; the code below was
re-executed to confirm. Following the six-step pattern:

\begin{verbatim}
from pulp import *

prob = LpProblem("Three_Var", LpMaximize)          # Step 1: problem
x1 = LpVariable("X1", lowBound=0)                  # Step 2: variables
x2 = LpVariable("X2", lowBound=0)
x3 = LpVariable("X3", lowBound=0)

prob += 4*x1 + 6*x2 + 5*x3, "Objective"            # Step 3: objective
prob += 2*x1 + x2 + x3 <= 9, "C1"                # Step 4: constraints
prob += x1 + 3*x2 + 2*x3 <= 14, "C2"
prob += x1 + x2 + 2*x3 <= 10, "C3"

prob.solve(PULP_CBC_CMD(msg=0))                    # Step 5: solve
print(LpStatus[prob.status], value(prob.objective))
for v in prob.variables():                         # Step 6: results
    print(v.name, "=", v.varValue)
\end{verbatim}

Actual output:

\begin{verbatim}
Optimal 35.0
X1 = 2.0
X2 = 2.0
X3 = 3.0
\end{verbatim}

The optimal objective value is $35$ at $(x_1, x_2, x_3) = (2, 2, 3)$.
All three constraints are binding ($4+2+3 = 9$, $2+6+6 = 14$, $2+2+6 = 10$);

\begin{qaflag}
Resolved 2026-07-10: the original data for this exercise matched a well-known copyrighted textbook example coefficient for coefficient. The data was replaced with an original instance (verified optimum $(2,2,3)$, value $35$) in both the book and this manual.
\end{qaflag}

% ----------------------------------------------------------------------------
\exsol{12.3}{New Demand Data}{1}

Only the \texttt{demand} dictionary changes.

\begin{verbatim}
from pulp import *

plants = ['Chicago', 'Denver']
markets = ['NYC', 'LA', 'Houston']

supply = {'Chicago': 100, 'Denver': 150}
demand = {'NYC': 100, 'LA': 80, 'Houston': 70}   # updated demand

cost = {
    ('Chicago', 'NYC'): 3, ('Chicago', 'LA'): 5, ('Chicago', 'Houston'): 2,
    ('Denver', 'NYC'): 4, ('Denver', 'LA'): 2, ('Denver', 'Houston'): 3
}

prob = LpProblem("Transportation", LpMinimize)
routes = [(p, m) for p in plants for m in markets]
x = LpVariable.dicts("Ship", routes, lowBound=0)
prob += lpSum([cost[r] * x[r] for r in routes]), "Total_Cost"
for p in plants:
    prob += lpSum([x[p,m] for m in markets]) <= supply[p], f"Supply_{p}"
for m in markets:
    prob += lpSum([x[p,m] for p in plants]) >= demand[m], f"Demand_{m}"

prob.solve(PULP_CBC_CMD(msg=0))
print(f"Status: {LpStatus[prob.status]}")
print(f"Total Cost: ${value(prob.objective)}")
for p in plants:
    for m in markets:
        if x[p,m].varValue > 0:
            print(f"  {p} -> {m}: {x[p,m].varValue}")
\end{verbatim}

Actual output:

\begin{verbatim}
Status: Optimal
Total Cost: $670.0
  Chicago -> NYC: 100.0
  Denver -> LA: 80.0
  Denver -> Houston: 70.0
\end{verbatim}

The new total cost is $\$670$ (up from $\$610$). CBC's plan sends all of
Chicago's supply to NYC and covers LA and Houston from Denver. (There are
alternative optima with the same cost, e.g., Denver serving LA and NYC while
Chicago covers Houston and part of NYC; students may report a different plan
at cost $670$.)

One-sentence explanation: total demand rose from $240$ to $250$ units and
shifted from LA (which Denver serves at the system's cheapest rate, $\$2$)
toward NYC (cheapest inbound rate $\$3$), so the network must ship more units
and ship them on more expensive routes. Note also that demand now exactly
equals total supply ($250 = 250$), so neither plant has spare capacity left to
exploit cheaper routings.

% ----------------------------------------------------------------------------
\exsol{12.4}{A Three-Plant Transportation Model}{2}

\begin{verbatim}
from pulp import *

# === DATA ===
plants = ['A', 'B', 'C']
markets = ['M1', 'M2', 'M3']
supply = {'A': 60, 'B': 70, 'C': 80}
demand = {'M1': 50, 'M2': 60, 'M3': 70}
cost = {
    ('A','M1'): 4, ('A','M2'): 6, ('A','M3'): 8,
    ('B','M1'): 5, ('B','M2'): 4, ('B','M3'): 3,
    ('C','M1'): 6, ('C','M2'): 7, ('C','M3'): 5
}

# === MODEL ===
prob = LpProblem("Three_Plant_Transport", LpMinimize)
routes = [(p, m) for p in plants for m in markets]
x = LpVariable.dicts("Ship", routes, lowBound=0)

prob += lpSum([cost[r] * x[r] for r in routes]), "Total_Cost"
for p in plants:
    prob += lpSum([x[p,m] for m in markets]) <= supply[p], f"Supply_{p}"
for m in markets:
    prob += lpSum([x[p,m] for p in plants]) >= demand[m], f"Demand_{m}"

# === SOLVE / RESULTS ===
prob.solve(PULP_CBC_CMD(msg=0))
print(f"Status: {LpStatus[prob.status]}")
print(f"Total Cost: {value(prob.objective)}")
for p in plants:
    shipped = sum(x[p,m].varValue for m in markets)
    print(f"Plant {p}: shipped {shipped} of {supply[p]}"
          f" (unused {supply[p]-shipped})")
for p in plants:
    for m in markets:
        if x[p,m].varValue > 0:
            print(f"  {p} -> {m}: {x[p,m].varValue}")
\end{verbatim}

Actual output:

\begin{verbatim}
Status: Optimal
Total Cost: 770.0
Plant A: shipped 50.0 of 60 (unused 10.0)
Plant B: shipped 70.0 of 70 (unused 0.0)
Plant C: shipped 60.0 of 80 (unused 20.0)
  A -> M1: 50.0
  B -> M2: 60.0
  B -> M3: 10.0
  C -> M3: 60.0
\end{verbatim}

Minimum total cost: $\$770$. Plan: A supplies all of M1 (50 units at cost 4),
B supplies all of M2 (60 at 4) plus 10 units of M3 (at 3), and C supplies the
remaining 60 units of M3 (at 5). Plants A and C have unused capacity (10 and
20 units respectively); B, whose routes to M2 and M3 are the cheapest in
their columns, is used to capacity. Sanity check by hand: the greedy
assignment M1$\leftarrow$A, M2$\leftarrow$B, M3$\leftarrow$B then C gives
$200 + 240 + 30 + 300 = 770$, and the LP confirms no cheaper plan exists.

% ----------------------------------------------------------------------------
\exsol{12.5}{Duals and Slacks from a Solved Model}{2}

Complete script (original product mix data, profit $3X_1 + 2X_2$):

\begin{verbatim}
from pulp import *

prob = LpProblem("Product_Mix", LpMaximize)
x1 = LpVariable("X1", lowBound=0)
x2 = LpVariable("X2", lowBound=0)
prob += 3*x1 + 2*x2, "Total_Profit"
prob += 10*x1 + 5*x2 <= 300, "Material"
prob += 4*x1 + 4*x2 <= 160, "Labor"
prob += 2*x1 + 6*x2 <= 180, "Machine"
prob.solve(PULP_CBC_CMD(msg=0))

print(LpStatus[prob.status], value(prob.objective), x1.varValue, x2.varValue)
for name, c in prob.constraints.items():
    print(name, "dual =", c.pi, "slack =", c.slack)
\end{verbatim}

Actual output:

\begin{verbatim}
Optimal 100.0 20.0 20.0
Material dual = 0.2 slack = -0.0
Labor dual = 0.25 slack = -0.0
Machine dual = -0.0 slack = 20.0
\end{verbatim}

\begin{enumerate}
\item Material: dual $0.2$, slack $0$. Labor: dual $0.25$, slack $0$.
Machine: dual $0$, slack $20$. (CBC prints $-0.0$; that is a signed floating
point zero, i.e., $0$.) Hand check of the duals: the binding rows give
$10y_1 + 4y_2 = 3$ and $5y_1 + 4y_2 = 2$, so $y_1 = 0.2$, $y_2 = 0.25$.
\item The Machine constraint has dual value $0$ and slack $20$: at the
optimum $(20,20)$, machine usage is $2(20) + 6(20) = 160 < 180$. A resource
with leftover capacity cannot be worth anything at the margin, so its shadow
price is zero. This is complementary slackness: for each constraint, the dual
value or the slack (or both) must be zero.
\item The Labor dual is $0.25$, so one extra hour of labor raises the optimal
profit by $\$0.25$, from $\$100$ to $\$100.25$ (valid as long as the current
basis stays optimal; Exercise 12.9 shows this rate holds up to $b = 168$).
\end{enumerate}

% ----------------------------------------------------------------------------
\exsol{12.6}{An Epsilon-Constraint Sweep}{2}

Following the template from the multi-objective chapter (Section 13.3), we
re-solve the product mix LP with the extra cap $2X_1 + 6X_2 \leq \varepsilon$.

\begin{verbatim}
from pulp import *

print(f"{'eps':>4} {'X1':>6} {'X2':>6} {'profit':>7}")
for eps in [60, 90, 120, 150, 180]:
    prob = LpProblem("Product_Mix_Eps", LpMaximize)
    x1 = LpVariable("X1", lowBound=0)
    x2 = LpVariable("X2", lowBound=0)
    prob += 3*x1 + 2*x2
    prob += 10*x1 + 5*x2 <= 300, "Material"
    prob += 4*x1 + 4*x2 <= 160, "Labor"
    prob += 2*x1 + 6*x2 <= 180, "Machine"
    prob += 2*x1 + 6*x2 <= eps, "Machine_cap"
    prob.solve(PULP_CBC_CMD(msg=0))
    print(f"{eps:>4} {x1.varValue:>6.1f} {x2.varValue:>6.1f}"
          f" {value(prob.objective):>7.1f}")
\end{verbatim}

Actual output:

\begin{verbatim}
 eps     X1     X2  profit
  60   30.0    0.0    90.0
  90   27.0    6.0    93.0
 120   24.0   12.0    96.0
 150   21.0   18.0    99.0
 180   20.0   20.0   100.0
\end{verbatim}

Trade-off: tightening the machine cap costs profit, but cheaply. From
$\varepsilon = 160$ (where the unconstrained optimum $(20,20)$ uses exactly
$160$ machine hours) down to $\varepsilon = 60$, the cap is binding and each
machine hour given up costs exactly $\$0.10$ of profit (the cap's shadow
price): profit falls linearly $100 \to 99 \to 96 \to 93 \to 90$ as
$\varepsilon$ falls $160 \to 150 \to 120 \to 90 \to 60$. For
$\varepsilon \geq 160$ (including the printed $\varepsilon = 180$) the cap is
slack and the solution is the original optimum. Along the sweep the plan
shifts steadily from the machine-heavy mix toward Product 1 only: the company
can cut machine usage by nearly two thirds (from $160$ to $60$ hours) while
losing only $10\%$ of profit. Each row is a point on the (profit, machine
usage) Pareto frontier.

% ----------------------------------------------------------------------------
\exsol{12.7}{The Constraint That Wasn't There}{2}

\begin{enumerate}
\item Prediction and experiment agree: PuLP neither raises an error nor
prints a warning; it happily solves the model you gave it. PuLP cannot know
that a constraint is missing; the two-constraint LP is a perfectly valid
problem. The run returns status \texttt{Optimal}.

\begin{verbatim}
from pulp import *

prob = LpProblem("Product_Mix_NoLabor", LpMaximize)
x1 = LpVariable("X1", lowBound=0)
x2 = LpVariable("X2", lowBound=0)
prob += 3*x1 + 2*x2, "Total_Profit"
prob += 10*x1 + 5*x2 <= 300, "Material"
prob += 2*x1 + 6*x2 <= 180, "Machine"     # Labor line is missing
prob.solve(PULP_CBC_CMD(msg=0))
print(LpStatus[prob.status], value(prob.objective), x1.varValue, x2.varValue)
print("Labor used:", 4*x1.varValue + 4*x2.varValue)
\end{verbatim}

Actual output:
\begin{verbatim}
Optimal 102.0 18.0 24.0
Labor used: 168.0
\end{verbatim}

\item The returned solution is $(18, 24)$ with profit $\$102$: the
intersection of the Material and Machine lines
($10x_1 + 5x_2 = 300$, $2x_1 + 6x_2 = 180$). The profit exceeds $\$100$
because dropping a constraint enlarges the feasible region, and the true
optimum $(20,20)$ had Labor binding, so relaxing Labor strictly helps. The
plan is \emph{not} implementable: it needs $4(18) + 4(24) = 168$ labor hours
but only $160$ are available.

\item With \texttt{prob += 4*x1 + 4*x2} (no \texttt{<= 160}), PuLP treats the
expression as a new \emph{objective} and replaces $3x_1 + 2x_2$ with
$4x_1 + 4x_2$. Recent PuLP versions print
\begin{verbatim}
UserWarning: Overwriting previously set objective.
\end{verbatim}
but the model still solves and reports \texttt{Optimal} (our run returns
$(18, 24)$ with objective value $168$, which is the value of $4x_1 + 4x_2$,
not profit). This bug is more dangerous than an omitted line because it does
two bad things at once, both invisibly: the labor constraint is still missing
\emph{and} the objective is now the wrong function, yet every status check
passes. An omitted line at least leaves the objective intact. Defensive
habits: print \texttt{prob} (or \texttt{prob.objective}) before solving, and
treat any \texttt{UserWarning} from PuLP as an error.
\end{enumerate}

\begin{qaflag}
Book fix applied (sync note): the warn box in the product mix example
(Subsection 12.12.3, Step 4) said PuLP
\emph{silently} overwrites the objective. Verified with PuLP 3.3.2 that a
\texttt{UserWarning} is printed. The book text was updated to say a warning
is printed but the model still solves. Exercise wording itself needed no
change; part (1) of this exercise refers to the omitted-line case, which is
truly silent.
\end{qaflag}

% ----------------------------------------------------------------------------
\exsol{12.8}{Minimize or Maximize?}{2}

\begin{enumerate}
\item Minimizing $3x_1 + 2x_2$ over the feasible region returns
$(x_1, x_2) = (0, 0)$ with objective $0$, status \texttt{Optimal}. Verified:

\begin{verbatim}
from pulp import *
prob = LpProblem("Product_Mix", LpMinimize)   # wrong sense
x1 = LpVariable("X1", lowBound=0)
x2 = LpVariable("X2", lowBound=0)
prob += 3*x1 + 2*x2, "Total_Profit"
prob += 10*x1 + 5*x2 <= 300, "Material"
prob += 4*x1 + 4*x2 <= 160, "Labor"
prob += 2*x1 + 6*x2 <= 180, "Machine"
prob.solve(PULP_CBC_CMD(msg=0))
print(LpStatus[prob.status], value(prob.objective), x1.varValue, x2.varValue)
\end{verbatim}

Output: \texttt{Optimal 0.0 0.0 0.0}. The model is not infeasible because
the constraints are all $\leq$ constraints with nonnegative data: the origin
satisfies every one of them. Minimizing a nonnegative objective over a region
containing the origin gives the useless but perfectly legal answer ``produce
nothing.''

\item Two one-line fixes, both verified to return the plan $(20, 20)$:
\begin{itemize}
\item Change the sense:
\texttt{prob = LpProblem("Product\_Mix", LpMaximize)}. \\
Output: \texttt{Optimal 100.0 20.0 20.0}.
\item Change the objective (keep \texttt{LpMinimize}): minimize the negated
profit, \texttt{prob += -3*x1 - 2*x2}. \\
Output: \texttt{Optimal -100.0 20.0 20.0}; the optimal plan is the same
$(20,20)$ and the profit is the negative of the reported objective,
$\$100$.
\end{itemize}

\item \texttt{LpStatus[prob.status]} only reports whether the solver finished
successfully on \emph{the model you built}; the wrong-sense model is a valid
LP with a legitimate optimum, so the status is \texttt{Optimal} in both the
buggy and the correct runs. Useful sanity checks: assert that the objective
value is strictly positive (a profit of \$0 from a company trying to make
money should raise eyebrows); assert the solution is not identically zero;
compare the objective against any known feasible plan (e.g., $(10,10)$ earns
\$50, so the optimum must be at least \$50); or check that at least one
resource constraint is binding, since a profit maximum with every resource
slack is suspicious.
\end{enumerate}

% ----------------------------------------------------------------------------
\exsol{12.9}{Profit as a Function of Labor}{3}

The book's Selected Solution is correct as printed; here is the full script
and its actual output (plans added for part 2's discussion).

\begin{verbatim}
from pulp import *

def solve_mix(b):
    prob = LpProblem("Product_Mix", LpMaximize)
    x1 = LpVariable("X1", lowBound=0)
    x2 = LpVariable("X2", lowBound=0)
    prob += 3*x1 + 2*x2
    prob += 10*x1 + 5*x2 <= 300, "Material"
    prob += 4*x1 + 4*x2 <= b, "Labor"
    prob += 2*x1 + 6*x2 <= 180, "Machine"
    prob.solve(PULP_CBC_CMD(msg=0))
    return value(prob.objective), prob.constraints["Labor"].pi, \
           x1.varValue, x2.varValue

for b in range(100, 200, 10):
    profit, dual, v1, v2 = solve_mix(b)
    print(f"b = {b}: profit = {profit:<6} labor dual = {dual:<5}"
          f" plan = ({v1}, {v2})")
\end{verbatim}

Actual output:

\begin{verbatim}
b = 100: profit = 75.0   labor dual = 0.75  plan = (25.0, 0.0)
b = 110: profit = 82.5   labor dual = 0.75  plan = (27.5, 0.0)
b = 120: profit = 90.0   labor dual = 0.75  plan = (30.0, 0.0)
b = 130: profit = 92.5   labor dual = 0.25  plan = (27.5, 5.0)
b = 140: profit = 95.0   labor dual = 0.25  plan = (25.0, 10.0)
b = 150: profit = 97.5   labor dual = 0.25  plan = (22.5, 15.0)
b = 160: profit = 100.0  labor dual = 0.25  plan = (20.0, 20.0)
b = 170: profit = 102.0  labor dual = -0.0  plan = (18.0, 24.0)
b = 180: profit = 102.0  labor dual = -0.0  plan = (18.0, 24.0)
b = 190: profit = 102.0  labor dual = -0.0  plan = (18.0, 24.0)
\end{verbatim}

\begin{enumerate}
\item The ten profits are $75, 82.5, 90, 92.5, 95, 97.5, 100, 102, 102, 102$.
The curve is piecewise linear and concave with slope changes at $b = 120$ and
at $b = 168$ (the latter lies between the printed points $160$ and $170$; it
is where the plan $(18,24)$, the intersection of Material and Machine, uses
exactly $4(18) + 4(24) = 168$ labor hours).
\item The slope of the profit curve at each $b$ equals the printed shadow
price \texttt{c.pi} of the Labor constraint. For $b \leq 120$ only Labor
limits production, the plan is $(b/4, 0)$, and each hour is worth $0.75$
(profit rises $7.5$ per $10$ hours: $75 \to 82.5 \to 90$). For
$120 \leq b \leq 168$, Material and Labor are both binding, the plan slides
along the Material line, and an extra hour is worth only $0.25$ (profit rises
$2.5$ per $10$ hours). The strictly decreasing sequence of shadow prices
$0.75 \to 0.25 \to 0$ is exactly the concavity of the curve.
\item For $b \geq 168$ the optimum is pinned at $(18, 24)$, where Material
and Machine are binding and Labor is slack; extra labor has shadow price $0$
and profit stays flat at $\$102$. Buying more of a non-binding resource is
worthless. This is the right-hand-side sensitivity analysis of Chapter 10
computed by brute force: the loop reproduces the shadow price and its
allowable range (the dual stays $0.25$ precisely for
$120 \leq b \leq 168$) without any tableau work.
\end{enumerate}
